\chapter{中心力与两体问题}

\section{把两个运动物体化为一个相对坐标}
孤立的两个粒子通过$V(|\bm r_1-\bm r_2|)$相互作用。
令$\bm r=\bm r_1-\bm r_2$，总质量$M=m_1+m_2$，约化质量$\mu=m_1m_2/M$，则
\begin{equation}
 T=\frac12M\norm{\dot{\bm R}}^2+\frac12\mu\norm{\dot{\bm r}}^2,\qquad
 M\ddot{\bm R}=0,\qquad \mu\ddot{\bm r}=-V'(r)\nvec r.
\end{equation}
所以相对运动等价于质量为$\mu$的粒子在固定中心势中运动。
当$m_2\gg m_1$时$\mu\simeq m_1$，才恢复重物近似固定的情形。
以下的$E$和$\ell$均指相对运动的能量和角动量大小。

中心力没有关于力心的力矩，因此$\bm\ell=\mu\bm r\times\dot{\bm r}$是固定矢量。
当$\ell\ne0$时，运动限制在垂直于$\bm\ell$的平面；$\ell=0$时为径向运动。
在运动平面选极坐标，方位角正向与$\bm\ell$遵循右手定则，
\begin{equation}
 \ell=\mu r^2\dot\phi,\qquad \dv{A}{t}=\frac12r^2\dot\phi=\frac\ell{2\mu}.
 \label{eq:cm-area}
\end{equation}
面积速度恒定来自中心力，对任意中心势都成立，不只适用于反平方力。

\section{有效势与圆轨道稳定性}
利用角动量守恒消去$\dot\phi$，能量写成
\begin{equation}
 E=\frac12\mu\dot r^2+V_{\rm eff}(r),\qquad
 V_{\rm eff}(r)=V(r)+\frac{\ell^2}{2\mu r^2}.
 \label{eq:cm-effective}
\end{equation}
最后一项描述固定角动量下逼近中心所需增加的横向动能。
它是径向降维后的有效势，不能在原三维受力图中再加一遍。
圆轨道半径$r_0$满足
\begin{equation}
 V'(r_0)=\frac{\ell^2}{\mu r_0^3},\qquad
 \omega_r^2=\frac1\mu\left[V''(r_0)+\frac{3V'(r_0)}{r_0}\right].
\end{equation}
第二式是在保持$\ell$不变时展开有效势得到的小径向振动频率。
$\omega_r^2>0$给出径向线性稳定；角频率为$\Omega^2=V'(r_0)/(\mu r_0)$。
小径向振动每完成一周，方位角约增加$2\pi\Omega/\omega_r$，可据此分析近心点方向。
频率不相等并不自动意味着轨道不闭合，下述各向同性谐振子就是反例。

\begin{example}{幂律吸引力的稳定圆轨道}{cm-power-law}
设径向力$F_r=-kr^s$，$k>0$。则
\begin{equation}
 \Omega^2=\frac{k}{\mu}r_0^{s-1},\qquad
 \omega_r^2=(s+3)\Omega^2.
\end{equation}
故$s>-3$时圆轨道对小径向扰动稳定；$s=-3$时二阶检验失效，不能据此宣布稳定。
反平方力$s=-2$给出$\omega_r=\Omega$；各向同性谐振子$s=1$给出$\omega_r=2\Omega$。
\end{example}

\section{轨道方程与Kepler椭圆}
当$\ell\ne0$时可用$\phi$代替时间，令$u=1/r$。
由$\dot r=-(\ell/\mu)\d u/\d\phi$以及径向Newton方程得到Binet方程
\begin{equation}
 \dv[2]{u}{\phi}+u=-\frac{\mu}{\ell^2u^2}F_r(1/u).
\end{equation}
对$V=-k/r$，$k>0$，解为
\begin{equation}
 r(\phi)=\frac{p}{1+e\cos(\phi-\phi_0)},\qquad
 p=\frac{\ell^2}{\mu k},\qquad
 e^2=1+\frac{2E\ell^2}{\mu k^2}.
 \label{eq:cm-conic}
\end{equation}
这里$e$是偏心率，$\phi_0$指定近心点方向。最后一式可由近心点$\dot r=0$代入能量得到。
当$\ell\ne0$时，$E<0$对应椭圆（$e<1$），$E=0$对应抛物线，$E>0$对应双曲线；
$e=0$是圆。$\ell=0$的径向坠落不应通过除以$\ell$的推导讨论。

椭圆半长轴$a$、半短轴$b$满足
\begin{equation}
 E=-\frac{k}{2a},\qquad p=a(1-e^2),\qquad b=a\sqrt{1-e^2}.
\end{equation}
用面积$\pi ab$除以式\eqref{eq:cm-area}的面积速度，得到周期
\begin{equation}
 \mathcal T=2\pi\sqrt{\frac{\mu a^3}{k}}.
\end{equation}
万有引力$k=Gm_1m_2$时，$\mathcal T^2=4\pi^2a^3/[G(m_1+m_2)]$。
此处$a$是相对轨道的半长轴；两个物体绕质心轨道的半长轴分别为$m_2a/M$与$m_1a/M$。
若需要位置随时间变化，引入偏近点角$\mathcal E$：
\begin{equation}
 r=a(1-e\cos\mathcal E),\qquad
 \mathcal E-e\sin\mathcal E=\sqrt{\frac{k}{\mu a^3}}(t-t_{\rm p}).
\end{equation}
第二式是Kepler方程，通常先数值求$\mathcal E$再恢复位置，不能把椭圆轨道误当成匀速转动。

\section{从轨道到散射截面}
考虑从无穷远入射的排斥势$V=k/r$，$k>0$。
入射相对速度$v_\infty$、碰撞参数$b$给出
\begin{equation}
 E=\frac12\mu v_\infty^2,\qquad \ell=\mu v_\infty b,\qquad
 b=\frac{k}{2E}\cot\frac\chi2.
\end{equation}
$\chi\in(0,\pi]$为质心系中相对速度的偏转角。
对轴对称、单调的$b(\chi)$，入射环带$2\pi b\,|\d b|$对应立体角$2\pi\sin\chi\,|\d\chi|$，故
\begin{equation}
 \dv{\sigma}{\Omega}=\frac{b}{\sin\chi}\left|\dv b\chi\right|
 =\left(\frac{k}{4E}\right)^2\frac1{\sin^4(\chi/2)}.
\end{equation}
这是经典Rutherford截面。若一个散射角对应多个碰撞参数，必须对各支求和；
实验室系角度通常还要作两体运动学变换。
该结果假定点粒子、无屏蔽Coulomb势和经典轨道；小角度的发散提示长程相互作用，
不是任何真实装置都具有无限计数率。

\section{Virial关系的适用条件}
定义$G=\bm r\cdot\bm p$，则$\dot G=2T+\bm r\cdot\bm F$。
若长时间平均存在且$[G(t)-G(0)]/t\to0$，就有
\begin{equation}
 2\ave T=\ave{rV'(r)}.
\end{equation}
对$V=\alpha r^n$，$2\ave T=n\ave V$，其中势能零点采用该齐次形式。
Kepler束缚轨道给出$2\ave T=-\ave V$。
上述条件可由有界的$G$保证；无界散射轨道一般不能直接套用束缚态的Virial结论。

\paragraph{复习与阅读}
两体问题先分离质心，再用角动量固定运动平面，最后用有效势分析径向运动。
需要轨道形状时用Binet方程，需要时间信息时再做积分。
延伸阅读：Taylor的两体与中心力问题；Goldstein的中心力运动。
